\chapter{Mixed Strategies}\label{ch:b2:mixed-strategies}

A tennis server chooses, on every first serve, whether to aim at the receiver's left or right; the
receiver must guess. If the server always aimed at the side where he wins more often, the receiver would
learn it and lean that way, and the advantage would vanish. Game theory predicts that a server who plays
well mixes his serves so that he wins equally often whichever side he picks. The chart tests that
prediction on 40 ``point games'' (one server, serving to one court, in one match) from ten professional
finals of 1974--1997.

\begin{openingchart}{Share of points won by the server after a first serve to the left and to the right,
  in 40 point games from ten professional tennis finals, 1974--1997. Each dot is one server on one court in
  one match; the line marks equal rates.}{Walker and Wooders (2001), Table~1}
\begin{tikzpicture}
\begin{axis}[omrecord, xmin=0.45, xmax=0.9, ymin=0.45, ymax=0.9, xlabel={Win rate after serving left},
    ylabel={Win rate after serving right}, width=0.55\linewidth, height=0.55\linewidth,
    xtick={0.5,0.6,0.7,0.8,0.9}, ytick={0.5,0.6,0.7,0.8,0.9}]
  \addplot[black!50, dashed, domain=0.45:0.9, samples=2, forget plot] {x};
  \addplot[only marks, mark=*, mark size=1.8pt, omDef] table[x=rate_l, y=rate_r]
    {figdata/bachelor-2/11-mixed-strategies/out/serves.csv};
\end{axis}
\end{tikzpicture}
\end{openingchart}

% ledger: b2:ww2001-t1, b2:ww2001-totals
% figdata: serves_summary (games 40, overall_l 0.641, overall_r 0.654, mean_abs_gap 0.085, p05 1, p10 2)
The dots lie around the line. Over all 3{,}026 first serves, servers won 64 per cent of points after serving
left and 65 per cent after serving right. In single point games the two rates differ, by 8.5 points on
average, but with 55 to 99 first serves per point game, split unevenly between the sides, such gaps are what chance
produces: a statistical test of equal
rates rejects in one point game of 40 at the 5 per cent level, and in two at the 10 per cent level, about as
many as chance alone would give.

\Cref{ch:b2:static-games} found games, such as matching pennies, with no \omterm{def:b2:static-games:nash-equilibrium}{Nash equilibrium} in which each
player picks one strategy for sure. This chapter allows players to randomise. It shows that every finite
game then has an equilibrium, how to compute it, and what it predicts: not which side a server picks on a
given point, but that his chances are equal on both sides, and that his choices follow no pattern the
receiver could exploit.

\begin{recall}
From \cref{ch:b2:static-games}: \omterm{def:b2:static-games:normal-form}{normal-form games}, \omterm{def:b2:static-games:best-response}{best responses}, \omterm{def:b2:static-games:nash-equilibrium}{Nash equilibrium}, and matching pennies.
From \cref{ch:b2:expected-utility}: expected payoffs over lotteries.
\end{recall}

\begin{notation}[Symbols of this chapter]\label{not:b2:mixed-strategies:symbols}
In $2 \times 2$ games, $p$ is the row player's probability of playing her first strategy and $q$ the column
player's probability of playing his; $v$ is the \omterm{def:b2:mixed-strategies:value}{value of a zero-sum game}. $\sigma_i$ is a \omterm{def:b2:mixed-strategies:mixed-strategy}{mixed strategy} of
player $i$. In the \omterm{def:b2:mixed-strategies:inspection-game}{inspection game}, $D$ is the tax due, $F$ the fine, $k$ the cost of an audit, $a$ the
probability of an audit and $e$ the probability of evasion.
\end{notation}

\section{Randomised strategies}\label{sec:b2:mixed-strategies:randomised}

\begin{definition}[Mixed strategy]\label{def:b2:mixed-strategies:mixed-strategy}
A \emph{mixed strategy}\index{mixed strategy} $\sigma_i$ of player $i$ is a probability distribution over her
strategies $S_i$, which in this context are called pure strategies. Players randomise independently, and
each evaluates a profile of mixed strategies by her expected payoff.
\end{definition}

\begin{definition}[Support of a mixed strategy]\label{def:b2:mixed-strategies:support}
The \emph{support}\index{support (of a mixed strategy)} of a \omterm{def:b2:mixed-strategies:mixed-strategy}{mixed strategy} is the set of pure strategies
to which it gives positive probability.
\end{definition}

A \omterm{def:b2:static-games:nash-equilibrium}{Nash equilibrium} in mixed strategies is a profile of mixed strategies, each a \omterm{def:b2:static-games:best-response}{best response} to the
others. In matching pennies, suppose the column player shows heads with probability $q$. The row player, who
wins when the coins match, earns $q - (1 - q) = 2q - 1$ from heads and $1 - 2q$ from tails. If $q > \tfrac12$ she
plays heads; if $q < \tfrac12$, tails; at $q = \tfrac12$ she is indifferent, and any mixture is a \omterm{def:b2:static-games:best-response}{best
response}. The same reasoning for the column player gives the only equilibrium: both play each side with
probability one half.

\section{The indifference condition}\label{sec:b2:mixed-strategies:indifference}

\begin{proposition}[Indifference]\label{prop:b2:mixed-strategies:indifference}
In a \omterm{def:b2:static-games:nash-equilibrium}{Nash equilibrium} in mixed strategies, every pure strategy in the \omterm{def:b2:mixed-strategies:support}{support} of a player's \omterm{def:b2:mixed-strategies:mixed-strategy}{mixed strategy}
gives her the same expected payoff, and no pure strategy outside the \omterm{def:b2:mixed-strategies:support}{support} gives her more.
\end{proposition}

\begin{proof}
A player's expected payoff from a \omterm{def:b2:mixed-strategies:mixed-strategy}{mixed strategy} is the average of the expected payoffs of the pure
strategies in its \omterm{def:b2:mixed-strategies:support}{support}, weighted by their probabilities. If one strategy in the \omterm{def:b2:mixed-strategies:support}{support} paid less than
another, moving its probability to the better one would raise the average, so the \omterm{def:b2:mixed-strategies:mixed-strategy}{mixed strategy} would not
be a \omterm{def:b2:static-games:best-response}{best response}. If a pure strategy outside the \omterm{def:b2:mixed-strategies:support}{support} paid more than the common payoff of those inside,
playing it for sure would be better.
\end{proof}

The proposition gives the method. To find a mixed equilibrium, choose each player's mixture so that the
\emph{other} player is indifferent among the strategies in his \omterm{def:b2:mixed-strategies:support}{support}. The striking consequence is that a
player's equilibrium mixture depends on the opponent's payoffs, not on her own.

Take a stylised penalty kick, used in laboratory experiments with professional footballers. The kicker
kicks left or right; the goalkeeper covers left or right. The kicker scores with probability 0.60 if both
choose left, 0.95 if he kicks left and the keeper covers right, 0.90 if he kicks right and the keeper covers
left, 0.70 if both choose right.
% ledger: b2:wooders2008-pk
The game has no pure equilibrium. If the keeper covers left with probability $q$, the kicker scores $0.60q +
0.95(1 - q)$ by kicking left and $0.90q + 0.70(1 - q)$ by kicking right; equal when $0.25 = 0.55q$, so $q =
0.4545$. If the kicker kicks left with probability $p$, the keeper concedes $0.60p + 0.90(1 - p)$ by covering
left and $0.95p + 0.70(1 - p)$ by covering right; equal when $0.20 = 0.55p$, so $p = 0.3636$. At these mixtures the
kicker scores with probability $0.60 \times 0.4545 + 0.95 \times 0.5455 = 0.7909$ whichever side he picks.

\begin{omfigure}
\begin{tikzpicture}[econ]
\begin{axis}[omrecord, grid=none, xmin=-0.04, xmax=1.04, ymin=-0.04, ymax=1.04,
    xlabel={$p$: kicker's probability of kicking left}, ylabel={$q$: keeper's probability of covering left},
    xtick={0,0.3636,1}, xticklabels={0, 0.36, 1}, ytick={0,0.4545,1}, yticklabels={0, 0.45, 1},
    width=0.55\linewidth, height=0.55\linewidth]
  \addplot[omDef, very thick] coordinates {(1,0) (1,0.4545) (0,0.4545) (0,1)};
  \addplot[omThm, very thick, dashed] coordinates {(0,0) (0.3636,0) (0.3636,1) (1,1)};
  \node[eqpoint] at (axis cs:0.3636,0.4545) {};
  \node[anchor=south west, font=\footnotesize] at (axis cs:0.38,0.47) {equilibrium};
  \node[anchor=west, font=\footnotesize, omDef, align=left] at (axis cs:0.55,0.3) {kicker's best\\response};
  \node[anchor=west, font=\footnotesize, omThm, align=left] at (axis cs:0.45,0.85) {keeper's best\\response};
\end{axis}
\end{tikzpicture}

{\small \omterm{def:b2:static-games:best-response}{Best responses} in the penalty-kick game. The kicker kicks left for sure when the keeper covers left
less than 45 per cent of the time, right when more, and is indifferent at 45 per cent. The keeper covers left
when the kicker kicks left more than 36 per cent of the time. The two best-response lines cross only at the
mixed equilibrium.}
\end{omfigure}

The kicker kicks right 64 per cent of the time and the keeper covers right 55 per cent of the time. Each
mixture is pinned down by the other player's indifference, so a change in any scoring probability moves both,
sometimes in surprising directions (\cref{exo:b2:mixed-strategies:6}).

\section{Existence}\label{sec:b2:mixed-strategies:existence}

\begin{theorem}[Existence of Nash equilibrium]\label{thm:b2:mixed-strategies:nash}
Every game with finitely many players, each with finitely many pure strategies, has at least one \omterm{def:b2:static-games:nash-equilibrium}{Nash
equilibrium} in mixed strategies (Nash 1950).
\end{theorem}

\admitted

The proof applies a fixed-point theorem to the players' best-response correspondences, and its idea can be
seen in the figure: in a $2 \times 2$ game each player's \omterm{def:b2:static-games:best-response}{best response} is a step that runs from one side of the
square to the other, so the two steps must cross (\cref{exo:b2:mixed-strategies:11}). The theorem says
nothing about how many equilibria there are or how players would find one.

\section{Zero-sum games and the minimax theorem}\label{sec:b2:mixed-strategies:minimax}

\begin{definition}[Zero-sum game]\label{def:b2:mixed-strategies:zero-sum}
A two-player game is \emph{zero-sum}\index{zero-sum game} if, in every profile, the two payoffs sum to zero:
one player's gain is the other's loss. A game in which the payoffs sum to the same constant in every profile,
such as the penalty kick, where the keeper's payoff is one minus the kicker's, is equivalent.
\end{definition}

\begin{definition}[Maximin strategy]\label{def:b2:mixed-strategies:maximin}
A \emph{maximin strategy}\index{maximin strategy} of a player is a \omterm{def:b2:mixed-strategies:mixed-strategy}{mixed strategy} that maximises the payoff
she can guarantee, the smallest expected payoff she could receive whatever the opponent does.
\end{definition}

\begin{definition}[Value of a zero-sum game]\label{def:b2:mixed-strategies:value}
The \emph{value}\index{value of a zero-sum game} $v$ of a \omterm{def:b2:mixed-strategies:zero-sum}{zero-sum game} is the payoff that the row player's
\omterm{def:b2:mixed-strategies:maximin}{maximin strategy} guarantees her when it equals the payoff that the column player's \omterm{def:b2:mixed-strategies:maximin}{maximin strategy} holds
her to: $\max_{\sigma_1}\min_{\sigma_2} u_1 = \min_{\sigma_2}\max_{\sigma_1} u_1 = v$.
\end{definition}

\begin{theorem}[Minimax theorem]\label{thm:b2:mixed-strategies:minimax}
In every finite two-player \omterm{def:b2:mixed-strategies:zero-sum}{zero-sum game}, $\max_{\sigma_1}\min_{\sigma_2} u_1(\sigma_1, \sigma_2) =
\min_{\sigma_2}\max_{\sigma_1} u_1(\sigma_1, \sigma_2)$. Maximin strategies form a \omterm{def:b2:static-games:nash-equilibrium}{Nash equilibrium}, and every
\omterm{def:b2:static-games:nash-equilibrium}{Nash equilibrium} consists of maximin strategies.
\end{theorem}

\begin{proof}[Proof for $2 \times 2$ games]
Always $\max_{\sigma_1}\min_{\sigma_2} u_1 \le \min_{\sigma_2}\max_{\sigma_1} u_1$: whatever player 1
guarantees with some $\sigma_1$, no $\sigma_2$ can hold her below it. If the game has a pure equilibrium (a saddle point), its payoff is the least in its row
and the greatest in its column, so the row player guarantees it and the column player holds her to it:
equality. Otherwise, by \cref{prop:b2:mixed-strategies:indifference} there are mixtures $p^*$ and $q^*$ that make
each player indifferent between his two strategies. Against $q^*$, every strategy of player 1 earns the same
amount $v$, so $\min\max \le v$; with $p^*$, player 1 earns $v$ against every strategy of player 2, so $\max\min \ge
v$. Hence $v \le \max\min \le \min\max \le v$. The general case is \admitted.
\end{proof}

In the penalty kick the \omterm{def:b2:mixed-strategies:value}{value} is 0.7909: the kicker can guarantee scoring with that probability, and the
keeper can hold him to it. Minimax play has a second implication, which tests often use: a player whose
choices followed a pattern, such as alternating, would be exploitable by an opponent who noticed it.
Equilibrium choices are serially independent.

\section{Tested on professional sport}\label{sec:b2:mixed-strategies:sport}

Professional sport offers \omterm{def:b2:mixed-strategies:zero-sum}{zero-sum games} played for high stakes by experienced players, with outcomes
that are recorded. The chart at the start of the chapter is one test; the listing shows what data from
equilibrium play would look like. It simulates 20 pairs of kicker and keeper playing 150 penalty kicks each
at the equilibrium mixtures, and tests, for each pair, whether the kicker's scoring rates on the two sides
differ.

\omcode{code/bachelor-2/11-mixed-strategies/listings/penalty_play.py}{1}{26}{Simulated equilibrium play of
the penalty-kick game, with a test of equal scoring rates for each
pair.}\label{lst:b2:mixed-strategies:penalty}

\omoutput{code/bachelor-2/11-mixed-strategies/outputs/penalty_play.txt}{Equilibrium play produces equal
scoring rates up to chance.}\label{out:b2:mixed-strategies:penalty}

% code: listings/penalty_play.py
Even under perfect equilibrium play, a 5 per cent test rejects equality in about one pair in twenty by
chance; in this run, none. The Wimbledon data passed the same kind of test: one rejection in 40 point games.

\begin{record}[Serves at Wimbledon]
% ledger: b2:ww2001-t1, b2:ww2001-totals
Walker and Wooders (2001) recorded the direction of every first serve, and whether the server won the
point, in ten finals between 1974 and 1997. Serves to the centre, 6
per cent of the total, were left out. Over the 40 point games, equality of the win rates after left and right
serves was rejected at the 5 per cent level once, and a joint test over all 40 gives a p-value of 0.852:
the win rates are consistent with equilibrium play. The same paper finds that servers switched direction
more often than independent choices would, so the serves were not serially independent; in laboratory data,
both equal win rates and serial independence were rejected.
\end{record}

\begin{model}{Mixed-strategy equilibrium}
\assumes Players who maximise expected payoffs, know the game, and correctly anticipate the opponent's
mixture.
\predicts Equal expected payoffs across the strategies each player uses; choices that are serially
independent; mixtures that respond to the opponent's payoffs, not one's own.
\tested Professional tennis serves (Walker and Wooders 2001): equal win rates hold, serial independence
does not; laboratory penalty kicks played by professionals and students (Palacios-Huerta and Volij 2008).
\breaks Serial correlation: servers switch too often; a re-examination of the laboratory data finds the
professionals' play inconsistent with minimax in several respects and the students' closer to it (Wooders
2008).
\end{model}

% ledger: b2:wooders2008-pk
The two findings are different tests of one model. Equal win rates are about what players achieve; serial
independence is about how they achieve it. Professionals pass the first test better than the second,
which is the more demanding: it requires a person to produce random sequences, which people do poorly.

\section{An inspection game}\label{sec:b2:mixed-strategies:inspection}

\begin{definition}[Inspection game]\label{def:b2:mixed-strategies:inspection-game}
An \emph{inspection game}\index{inspection game} is a game between an inspector, who decides whether to
check, at a cost, and an inspectee, who decides whether to comply or to violate a rule, gaining if she
violates unchecked and paying a penalty if she violates and is checked.
\end{definition}

A taxpayer owes $D$ and can declare it (honest) or not (evade). A tax inspector can audit at a cost $k$. An audit
of an evader recovers the tax and imposes a fine $F$; the inspector values the tax he collects, net of audit
costs, while fines go to the treasury. The taxpayer pays $D$ if honest, $D + F$ if she evades and is audited,
nothing if she evades unaudited. The inspector collects $D - k$ if he audits (from an honest taxpayer or a
caught evader), $D$ if he does not audit an honest taxpayer, and nothing from an unaudited evader.

\begin{omfigure}
\begin{tikzpicture}[econ, font=\footnotesize]
  \foreach \ox/\title/\sa/\sb/\ta/\tb/\p/\q/\r/\s in {
      0/Penalty kick (scoring probability)/Left/Right/Left/Right/{0.60}/{0.95}/{0.90}/{0.70},
      7.0/Inspection game/Honest/Evade/Audit/No audit/{$-D$, $D - k$}/{$-D$, $D$}/{$-D - F$, $D - k$}/{0, 0}} {
    \node[font=\small\bfseries] at (\ox+2.3,1.9) {\title};
    \node[payoff, minimum width=25mm] at (\ox+1.0,0.4) {\p};
    \node[payoff, minimum width=25mm] at (\ox+3.6,0.4) {\q};
    \node[payoff, minimum width=25mm] at (\ox+1.0,-0.4) {\r};
    \node[payoff, minimum width=25mm] at (\ox+3.6,-0.4) {\s};
    \node[anchor=east] at (\ox-0.4,0.4) {\sa};
    \node[anchor=east] at (\ox-0.4,-0.4) {\sb};
    \node at (\ox+1.0,1.1) {\ta};
    \node at (\ox+3.6,1.1) {\tb};
  }
\end{tikzpicture}

{\small Two games without pure equilibria. Left: the kicker's scoring probability, rows for the kicker and
columns for the keeper; the keeper's payoff is one minus it. Right: the taxpayer (rows) and the inspector
(columns), the taxpayer's payoff first.}
\end{omfigure}

There is no pure equilibrium: if the inspector audits, the taxpayer is honest; then the inspector prefers
not to audit; then she evades; then he audits. In the mixed equilibrium each player's mixture makes the
other indifferent.

\begin{proposition}[Equilibrium of the inspection game]\label{prop:b2:mixed-strategies:inspection}
If $0 < k < D$, the \omterm{def:b2:mixed-strategies:inspection-game}{inspection game} has a unique equilibrium, in which the inspector audits with probability
$a^* = D/(D + F)$ and the taxpayer evades with probability $e^* = k/D$. A higher fine lowers the audit rate and
leaves the evasion rate unchanged.
\end{proposition}

\begin{proof}
The taxpayer is indifferent when $D = a(D + F)$, so $a^* = D/(D + F)$. The inspector earns $D - k$ from auditing
whatever the taxpayer does, and $(1 - e)D$ from not auditing; he is indifferent when $e^* = k/D$. The fine enters
only the taxpayer's indifference condition, which fixes the inspector's mixture.
\end{proof}

The result is a direct consequence of the indifference condition, and it runs against intuition: in
equilibrium, a harsher penalty does not deter evasion, because the inspector, anticipating more honesty,
audits less until the taxpayer is again indifferent. Evasion falls only if the cost of an audit falls or the
inspector's gain from catching an evader rises. The prediction assumes that the inspector chooses his audit
rate after seeing nothing and cannot commit to it; an inspector who can announce, and keep to, an audit rate
changes the game (\cref{pb:b2:mixed-strategies:1}, Part IV), a theme of \cref{ch:b2:dynamic-games}.

\section{Exercises}\label{sec:b2:mixed-strategies:exercises}

\begin{exercise}[$\star$]\label{exo:b2:mixed-strategies:1}
In a \omterm{def:b2:static-games:coordination-game}{coordination game} the row player gets 2 and the column player 1 if both choose $A$, the row player 1 and the
column player 2 if both choose $B$, and both get 0 otherwise. Find the pure equilibria and the mixed
equilibrium, and each player's expected payoff in the mixed one.
\end{exercise}

\begin{exercise}[$\star$]\label{exo:b2:mixed-strategies:2}
A server wins the point with probability 0.58 if he serves left and the receiver leans left, 0.79 if he serves
left and the receiver leans right, 0.73 if he serves right and the receiver leans left, and 0.49 if both choose
right. Find the equilibrium mixtures and the server's winning probability.
\end{exercise}

\begin{exercise}[$\star$]\label{exo:b2:mixed-strategies:3}
Read the opening chart. What were the overall win rates after serves to the left and to the right? In how many
point games was equality rejected at the 5 per cent level, and how many would you expect by chance?
\end{exercise}

\begin{exercise}[$\star$]\label{exo:b2:mixed-strategies:4}
What data from tennis matches would refute the prediction of equal win rates? Distinguish evidence against the
model from evidence of skill differences between players.
\end{exercise}

\begin{exercise}[$\star\star$]\label{exo:b2:mixed-strategies:5}
A \omterm{def:b2:mixed-strategies:zero-sum}{zero-sum game} has row payoffs 4 and 2 in the first row, 3 and 1 in the second. Find each player's \omterm{def:b2:mixed-strategies:maximin}{maximin
strategy} and the \omterm{def:b2:mixed-strategies:value}{value}. Why is no mixing needed?
\end{exercise}

\begin{exercise}[$\star\star$]\label{exo:b2:mixed-strategies:6}
Predict, then compute: the keeper in the penalty-kick game becomes better at saving kicks to his left when he
covers left, so that the scoring probability there falls from 0.60 to 0.50. How do the kicker's and the keeper's
mixtures change, and the \omterm{def:b2:mixed-strategies:value}{value}?
\end{exercise}

\begin{exercise}[$\star\star$]\label{exo:b2:mixed-strategies:7}
In one point game of the Wimbledon data, a server served left 33 times and won 20 points, and served right 26
times and won 22. Compute the two win rates, the pooled rate and the Pearson statistic $\sum (g - np)^2/(np(1 -
p))$ over the two directions. Is equality rejected at the 5 per cent level (critical value 3.84)?
\end{exercise}

\begin{exercise}[$\star\star$]\label{exo:b2:mixed-strategies:8}
In rock--paper--scissors each player wins 1 from the other when her choice beats his and nothing changes hands
on a tie. Use the indifference condition to find the equilibrium.
\end{exercise}

\begin{exercise}[$\star\star$]\label{exo:b2:mixed-strategies:9}
Design a test of serial independence for a tennis server's sequence of serve directions. What pattern would
reject it, and why might a player who knows the game still produce such a pattern?
\end{exercise}

\begin{exercise}[$\star\star\star$]\label{exo:b2:mixed-strategies:10}
For a $2 \times 2$ \omterm{def:b2:mixed-strategies:zero-sum}{zero-sum game} without a saddle point, with row payoffs $a_{11}, a_{12}, a_{21}, a_{22}$, derive
$p^*$, $q^*$ and the \omterm{def:b2:mixed-strategies:value}{value} $v = (a_{11}a_{22} - a_{12}a_{21})/(a_{11} - a_{12} - a_{21} + a_{22})$, and complete the proof
of \cref{thm:b2:mixed-strategies:minimax}.
\end{exercise}

\begin{exercise}[$\star\star\star$]\label{exo:b2:mixed-strategies:11}
Prove that every $2 \times 2$ game has a \omterm{def:b2:static-games:nash-equilibrium}{Nash equilibrium}, possibly in mixed strategies, by considering the
best-response correspondences in the unit square.
\end{exercise}

\begin{exercise}[$\star\star\star$]\label{exo:b2:mixed-strategies:12}
In the \omterm{def:b2:mixed-strategies:inspection-game}{inspection game}, suppose the inspector also receives a share $s$ of the fine. Find the equilibrium. Which
parameters now affect evasion, and in which direction?
\end{exercise}

\section{Problem: Brenland's tax inspectors}\label{sec:b2:mixed-strategies:problem}

\begin{problem}[{Weekend problem --- the tax office of Brenland, an invented economy, decides how often to
audit a group of self-employed taxpayers: how often do they evade, what does a higher fine achieve, and when
does it pay the office to announce its audit rate in advance}]
\label{pb:b2:mixed-strategies:1}
Each taxpayer owes $D = 1{,}000$ dalens. An audit costs the office $k = 200$ dalens. An evader who is audited pays
the tax and a fine $F = 1{,}500$; fines go to the treasury. The office maximises the tax it collects net of audit
costs.

\medskip
\noindent\textbf{Part I --- The inspection game.}
\begin{enumerate}
  \item Write the payoff table.
  \item Show that there is no pure-strategy equilibrium.
  \item Which player's payoffs determine the audit rate in a mixed equilibrium?
  \item Which player's payoffs determine the evasion rate?
  \item Does the taxpayer have a dominant strategy?
  \item Is the game \omterm{def:b2:mixed-strategies:zero-sum}{zero-sum}?
\end{enumerate}

\medskip
\noindent\textbf{Part II --- The mixed equilibrium.}
\begin{enumerate}[resume]
  \item Find the equilibrium audit probability.
  \item Find the equilibrium evasion probability.
  \item Find the taxpayer's expected payment.
  \item Find the office's expected net collection per taxpayer.
  \item What share of the tax due does the office collect, before audit costs?
  \item Check that at these probabilities neither player gains by changing alone.
\end{enumerate}

\medskip
\noindent\textbf{Part III --- A higher fine.}
\begin{enumerate}[resume]
  \item Predict what a fine of 4{,}000 dalens does to evasion and to audits.
  \item Compute the new equilibrium.
  \item What happens to the office's audit costs?
  \item What change would reduce evasion in this model?
  \item Why might real taxpayers respond to the fine even if the model says they do not?
  \item What evidence would distinguish the two views?
\end{enumerate}

\medskip
\noindent\textbf{Part IV --- Commitment.} The office can announce an audit rate and keep to it, at a fixed cost of
100 dalens per taxpayer (staff hired in advance). A taxpayer who is indifferent declares honestly.
\begin{enumerate}[resume]
  \item What is the lowest audit rate that makes every taxpayer honest?
  \item Find the office's net collection under commitment, with $F = 1{,}500$.
  \item Compare with Part II. Does commitment pay?
  \item Repeat with $F = 500$.
  \item Find the fine above which commitment pays.
  \item State the result: the equilibrium audit and evasion rates, and the fine at which commitment pays.
\end{enumerate}
\end{problem}
